% Figura: capas del modelo de datos. El crudo nunca se toca; cada capa agrega.
\begin{tikzpicture}[node distance=7mm]

\node[almacen,text width=46mm] (crudo) {\textbf{Capa 0 --- crudo}\\
  \cd{monitoreo\_sc\_electrico}\\ \cd{radiacion\_sc\_15s}\\
  \nota{tal como salió del CSV}};

\node[comp,text width=46mm,right=14mm of crudo] (corr) {\textbf{Capa 1 --- corrección}\\
  \cd{v\_sc\_electrico\_corregido}\\ \cd{v\_sc\_radiacion\_corregida}\\
  \nota{85\,°C, rangos, offset, pre-2025}};

\node[comp,text width=46mm,right=14mm of corr] (cal) {\textbf{Capa 2 --- calibración}\\
  \cd{v\_sc\_radiacion\_calibrada}\\
  \nota{W/m$^2$, kt*, \cd{qc\_ok}}};

\node[comp,text width=46mm,below=12mm of cal] (perf) {\textbf{Capa 3 --- desempeño}\\
  \cd{v\_sc\_performance}\\
  \nota{\cd{pr\_pv1}, \cd{pr\_pv2}}};

\node[comp,text width=42mm,below=12mm of corr] (cs) {\cd{radiacion\_sc\_clearsky}\\ \cd{radiacion\_sc\_poa}\\
  \nota{calculadas con pvlib}};

\draw[flecha] (crudo) -- node[etiq,above]{vista} (corr);
\draw[flecha] (corr) -- node[etiq,above]{vista} (cal);
\draw[flecha] (cs) -- (cal);
\draw[flecha] (cs) -- (perf);
\draw[flecha] (corr.south) |- (perf.west);

\node[etiq,below=3mm of crudo,align=center,fill=none]
  {sin pérdida de información:\\ nada se sobrescribe};

\end{tikzpicture}
